%% =====================================================================
%%  Unit 11 — Can I have a coffee, please?  (Au café et au restaurant)
%%  Pistes : 11.01–11.04 dialogues · 11.05–11.06 vocabulaire · 11.07–11.08 grammaire
%%           11.09 prononciation · 11.10–11.13 exercices (textes en annexe)
%% =====================================================================

\newcommand{\otIllusUnXI}{\otIllusDuo{latte/104/86/15/0.10,menu/190/92/20/0.09,fries/122/14/13/0.08,
  to go/150/8/11/0.08,check/196/22/12/0.07,tip/100/50/12/0.07}{sam}{What can I get you?}{lea}{A latte, please!}}

% Menu du restaurant (exercice 1)
\newcommand{\otMenu}{%
  \begin{tikzpicture}
    \fill[otOrangeTint,rounded corners=2mm] (0,0) rectangle (92mm,64mm);
    \draw[otOrangeDark,line width=0.6pt,rounded corners=1.6mm] (2mm,2mm) rectangle (90mm,62mm);
    \node[font=\otBold\fontsize{12}{14}\selectfont,text=otOrangeDark] at (46mm,55mm) {\enb{Brooklyn Bistro}};
    \node[font=\otMedium\fontsize{7}{8}\selectfont,text=otMuted] at (46mm,49.5mm) {\otCaps[12]{Menu}};
    \foreach \plat/\prix [count=\i] in {Soup of the day/\$6,Chicken salad/\$14,Burger and fries/\$16,
        Fish of the day/\$19,Ice cream/\$5,Coffee or tea/\$3,Water/free}{%
      \pgfmathsetmacro{\yy}{45-\i*5.6}
      \node[anchor=west,font=\fontsize{10}{12}\selectfont] at (8mm,\yy mm) {\enb{\plat}};
      \draw[otOrangeDark!40,line width=0.4pt,dash pattern=on 0.8pt off 1.6pt] (44mm,\yy mm-1mm) -- (74mm,\yy mm-1mm);
      \node[anchor=east,font=\otSemi\fontsize{10}{12}\selectfont] at (84mm,\yy mm) {\enb{\prix}};}
  \end{tikzpicture}}

\unite[Can I have a\\coffee, please?]{11}{Can I have a coffee, please?}{Au café et au restaurant}{otIllusUnXI}

\situation{Épisode 11}{Jeudi matin, Léa passe en coup de vent au café où travaille Sam. Le vendredi
soir, Jack invite Emma et Léa dans un petit restaurant de Brooklyn pour fêter la fin d'un gros
projet. Au moment de l'addition, Léa découvre une coutume américaine…}

\begin{objectifs}
\item commander au café et au restaurant ;
\item demander poliment : \en{Can I have…?}, \en{Could I have…?}, \en{I'd like…} ;
\item proposer quelque chose et répondre : \en{Would you like…? — Yes, please.} ;
\item demander l'addition et comprendre le système du pourboire ;
\item prononcer \en{can} sous sa forme faible \phon{kən}.
\end{objectifs}

\auprogramme
  {Commander, demander poliment, payer}
  {\en{Can I…?}, \en{Could I…?}\par\vspace{0.6mm}\en{I'd like}, \en{Would you like…?}}
  {Le menu, les plats, les boissons : 30 expressions}
  {Les formes faibles de \en{can} et \en{could}}
  {Le pourboire (\en{tipping})}

\semaine{casque/otNavy/Dialogues/Écouter et lire les deux scènes,
  bulle/otGreenDark/Vocabulaire/Apprendre les 30 expressions,
  blocs/otViolet/Grammaire/Demander poliment,
  sons/otOrangeDark/Sons et culture/can faible et le pourboire,
  crayon/otNavy/Exercices/12 exercices corrigés,
  coche/otGreenDark/Bilan/L'essentiel et le test express}

\clearpage
%% ---------------------------------------------------------------------
\rubrique{ecouter}{Commander}

\begin{dialogue}{1}{For here or to go?}{Jeudi, 8 h 45. Léa entre dans le café où Sam est barista.
  Elle est pressée.}{11.01}{11.02}
\dit{sam}{Hi, Léa! \hl{What can I get you}?\note{1}}{Salut, Léa ! Qu'est-ce que je te sers ?}
\dit{lea}{\hl{Can I have} a large latte, please?}{Je peux avoir un grand latte, s'il te plaît ?}
\dit{sam}{Sure. \hl{For here or to go}?\note{2}}{Bien sûr. Sur place ou à emporter ?}
\dit{lea}{To go, please. I'm late!}{À emporter, s'il te plaît. Je suis en retard !}
\dit{sam}{\hl{Would you like} a muffin with that?\note{3}}{Tu veux un muffin avec ça ?}
\dit{lea}{No, thanks. \hl{I'd like} a croissant, please.}{Non, merci. Je voudrais un croissant, s'il te plaît.}
\dit{sam}{Sorry, we don't have any croissants today.}{Désolé, on n'a pas de croissants aujourd'hui.}
\dit{lea}{OK, a muffin then! How much is that?}{Bon, alors un muffin ! Ça fait combien ?}
\dit{sam}{Seven fifty.}{Sept dollars cinquante.}
\dit{lea}{\hl{Here you are}.}{Tiens.}
\dit{sam}{Thanks! Have a great day!}{Merci ! Bonne journée !}
\end{dialogue}

\begin{comprendre}
\pointnote{1} \en{\textbf{What can I get you?}} « Qu'est-ce que je vous sers ? », la question rituelle
  du serveur. Pour commander : \en{Can I have…?} (« Est-ce que je peux avoir… ? »).
\pointnote{2} \en{\textbf{For here or to go?}} « Sur place ou à emporter ? ». Au Royaume-Uni, on dit
  plutôt \en{Eat in or take away?}
\pointnote{3} \en{\textbf{Would you like…?}} « Voulez-vous… ? », « désirez-vous… ? » : la façon polie
  de proposer. \en{I'd like} (= \en{I would like}) veut dire « je voudrais » : c'est la façon polie de
  demander.
\end{comprendre}

\begin{dialogue}{2}{Can we have the check, please?}{Vendredi, 20 heures, dans un restaurant de
  Brooklyn. Le serveur arrive à la table de Jack, Emma et Léa.}{11.03}{11.04}
\dit{serveur}{Good evening! \hl{Are you ready to order}?}{Bonsoir ! Vous avez choisi ?}
\dit{emma}{Yes. \hl{I'd like} the chicken salad, please.}{Oui. Je voudrais la salade au poulet, s'il vous plaît.}
\dit{jack}{And \hl{I'll have} the burger, with fries.\note{4}}{Et moi, je prendrai le burger, avec des frites.}
\dit{serveur}{And for you, ma'am?}{Et pour vous, madame ?}
\dit{lea}{\hl{Could I have} the fish, please?\note{5}}{Pourrais-je avoir le poisson, s'il vous plaît ?}
\dit{serveur}{Of course. \hl{Would you like} something to drink?}{Bien sûr. Vous désirez quelque chose à boire ?}
\dit{jack}{Three glasses of water, please.}{Trois verres d'eau, s'il vous plaît.}
\didascalie{Une heure plus tard.}
\dit{emma}{Excuse me! \hl{Can we have the check}, please?}{Excusez-moi ! L'addition, s'il vous plaît ?}
\dit{serveur}{Sure. Here you are.}{Bien sûr. La voici.}
\dit{lea}{Is \hl{the tip} included?\note{6}}{Le pourboire est compris ?}
\dit{jack}{No, it isn't. In the US, we add fifteen or twenty percent.}{Non. Aux États-Unis, on ajoute quinze ou vingt pour cent.}
\end{dialogue}

\begin{comprendre}
\pointnote{4} \en{\textbf{I'll have the burger.}} « Je prendrai le burger » : au restaurant, on commande
  souvent avec \en{I'll have…} (= \en{I will have}). \en{fries} = « des frites ».
\pointnote{5} \en{\textbf{Could I have…?}} « Pourrais-je avoir… ? » : \en{could} est encore plus poli
  que \en{can}. \en{ma'am} (\en{madam}) : « madame », pour s'adresser poliment à une cliente.
\pointnote{6} \en{\textbf{the check}, \textbf{the tip}} « l'addition », « le pourboire ». Au Royaume-Uni,
  l'addition se dit \en{the bill}. Voir la rubrique Culture.
\end{comprendre}

%% ---------------------------------------------------------------------
\rubrique{mots}{Les mots de l'unité}

Voici les 30 expressions de l'unité. Au café comme au restaurant, \en{please} et \en{thank you}
sont indispensables : en anglais, une demande sans \en{please} paraît vite impolie.

\begin{vocabulaire}{Au menu}{11.05}
\categorie{À boire}
\mot{a coffee}{ə ˈkɑːfi}{un café}
\mot{a latte}{ə ˈlɑːteɪ}{un café au lait}
\mot{a tea}{ə ˈtiː}{un thé}
\mot{a glass of water}{ə ˌɡlæs əv ˈwɑːtɚ}{un verre d'eau}
\mot{orange juice}{ˈɔːrɪndʒ ˌdʒuːs}{du jus d'orange}
\categorie{À manger}
\mot{the menu}{ðə ˈmenjuː}{la carte, le menu}
\mot{a muffin}{ə ˈmʌfɪn}{un muffin}
\mot{soup}{suːp}{de la soupe}
\mot{a salad}{ə ˈsæləd}{une salade}
\mot{a burger}{ə ˈbɝːɡɚ}{un hamburger}
\mot{fries}{fraɪz}{des frites}
\mot{a steak}{ə ˈsteɪk}{un steak}
\mot{dessert}{dɪˈzɝːt}{le dessert}
\mot{ice cream}{ˈaɪs ˌkriːm}{de la glace}
\categorie{Payer}
\mot{the check}{ðə ˈtʃek}{l'addition \registre{américain}}
\mot{a tip}{ə ˈtɪp}{un pourboire}
\end{vocabulaire}

\begin{vocabulaire}{Commander poliment}{11.06}
\categorie{Le serveur}
\mot{What can I get you?}{ˌwʌt kən aɪ ˈɡet juː}{Qu'est-ce que je vous sers ?}
\mot{Are you ready to order?}{ɑːr juː ˌredi tu ˈɔːrdɚ}{Vous avez choisi ?}
\mot{For here or to go?}{fɚ ˈhɪr ɔːr tə ˈɡoʊ}{Sur place ou à emporter ?}
\mot{Would you like something to drink?}{wʊd juː ˌlaɪk ˌsʌmθɪŋ tə ˈdrɪŋk}{Vous désirez quelque chose à boire ?}
\mot{Here you are.}{ˌhɪr juː ˈɑːr}{Voici. Tenez.}
\categorie{Le client}
\mot{Can I have a coffee, please?}{kən aɪ ˌhæv ə ˈkɑːfi pliːz}{Je peux avoir un café, s'il vous plaît ?}
\mot{Could I have the fish?}{kʊd aɪ ˌhæv ðə ˈfɪʃ}{Pourrais-je avoir le poisson ?}
\mot{I'd like a salad.}{aɪd ˌlaɪk ə ˈsæləd}{Je voudrais une salade.}
\mot{I'll have the burger.}{aɪl ˌhæv ðə ˈbɝːɡɚ}{Je prendrai le burger.}
\mot{Yes, please. · No, thank you.}{ˌjes ˈpliːz · ˌnoʊ ˈθæŋk juː}{Oui, volontiers. · Non, merci.}
\mot{To go, please.}{tə ˈɡoʊ pliːz}{À emporter, s'il vous plaît.}
\mot{Can we have the check, please?}{kən wiː ˌhæv ðə ˈtʃek pliːz}{L'addition, s'il vous plaît.}
\mot{Is the tip included?}{ɪz ðə ˌtɪp ɪnˈkluːdɪd}{Le pourboire est compris ?}
\mot{It's delicious!}{ɪts dɪˈlɪʃəs}{C'est délicieux !}
\end{vocabulaire}

%% ---------------------------------------------------------------------
\rubrique{grammaire}{Demander et proposer poliment}

\otIntertitre[otViolet]{1 · Demander : du plus simple au plus poli}
\begin{center}\otTableau
\begin{tabular}{@{}>{\enOn\bfseries}l<{\enOff}@{\hspace{8mm}}l@{\hspace{8mm}}l@{}}
\otTH{Formule} & \otTH{Sens} & \otTH{Registre}\\\midrule
Can I have a coffee, please?     & Je peux avoir un café ?       & courant\\
I'd like a coffee, please.       & Je voudrais un café.          & poli\\
I'll have a coffee, please.      & Je prendrai un café.          & au restaurant\\
Could I have a coffee, please?   & Pourrais-je avoir un café ?   & très poli\\\bottomrule
\end{tabular}\end{center}

\blocaudio{11.07}{Écoutez et répétez}{Les quatre formules du tableau, avec \en{please}, puis deux
demandes au pluriel avec \en{Can we…?}}

\otIntertitre[otViolet]{2 · Proposer : \en{Would you like…?}}
\begin{paire}
\begin{regle}[Proposer]
\exen{Would you like a muffin?}{Voulez-vous un muffin ?}
\exen{Would you like some tea?}{Voulez-vous du thé ?}
\exen{Would you like something to drink?}{Quelque chose à boire ?}
\end{regle}
\begin{regle}[Répondre]
\exen{Yes, please.}{Oui, volontiers.}
\exen{No, thank you. / No, thanks.}{Non, merci.}
\exen{Yes, I'd love some!}{Oui, avec plaisir !}
\end{regle}
\end{paire}

\blocaudio{11.08}{Écoutez et répétez}{Une proposition avec \en{Would you like…?}, puis la réponse.}

\begin{regle}[\en{I like} ou \en{I'd like} ?]
\en{I like coffee} = « j'aime le café » (en général, toujours). \en{I'd like a coffee} = « je
voudrais un café » (maintenant, c'est une demande). Ne les confondez pas : au café, \en{I like a
coffee} n'a pas de sens !
\end{regle}

\begin{attention}[Les pièges des francophones]
\fauxjuste{I want a coffee.}{I'd like a coffee, please.}
\fauxjuste{Give me the bill.}{Can we have the check, please?}
\fauxjuste{Would you like a tea? — Yes.}{Would you like a tea? — Yes, please.}
\fauxjuste{Can I have a water?}{Can I have a glass of water?}
\par\vspace{1mm}{\fontsize{10}{12.6}\selectfont \en{I want…} est grammaticalement juste, mais il
sonne sec, voire impoli, dans la bouche d'un client. Préférez \en{I'd like…} ou \en{Can I have…?}\par}
\end{attention}

%% ---------------------------------------------------------------------
\rubrique[Les formes faibles de can]{sons}{Les formes faibles de \en{can} et \en{could}}

Dans une demande, \en{can}, \en{could} et \en{would} ne sont pas accentués : leur voyelle s'affaiblit
en \phon{ə}. \en{Can I have…?} se dit \phon{kənaɪ ˈhæv}, presque « k'naï have ». En revanche, en fin
de phrase (\en{Yes, you can.}), \en{can} retrouve sa forme pleine \phon{kæn}.

\begin{phrases}{11.09}{Écoutez et répétez à vitesse normale, en n'accentuant que les mots importants.}
\phrase{Can I have a coffee?}{\phon{kən} · \textbf{coffee}}
\phrase{Can we have the check, please?}{\phon{kən} · \textbf{check}}
\phrase{Could I have the fish?}{\phon{kəd} · \textbf{fish}}
\phrase{Would you like some tea?}{\phon{wədʒə} · \textbf{like}, \textbf{tea}}
\phrase{I'd like a salad.}{\phon{aɪd} · \textbf{like}, \textbf{salad}}
\phrase{I'll have the burger.}{\phon{aɪl} · \textbf{have}, \textbf{burger}}
\phrase{Yes, you can.}{en fin de phrase, forme pleine : \phon{kæn}}
\end{phrases}

%% ---------------------------------------------------------------------
\rubrique{culture}{Le pourboire}

\begin{culture}[\en{Tipping} : une règle, pas un bonus]
Aux États-Unis, le service n'est \textbf{pas compris} : les serveurs vivent en grande partie de leurs
pourboires. Au restaurant, on laisse entre \textbf{15 et 20 \%} du montant (hors taxe), parfois plus si
le service est excellent. Ne pas laisser de pourboire est perçu comme une remarque très négative.
\end{culture}

\begin{paire}
\begin{culture}[Ailleurs]
Au comptoir d'un café, une boîte à pourboires (\en{tip jar}) ou l'écran de paiement propose souvent
d'ajouter 15, 18 ou 20 \% : c'est facultatif. Dans un bar, on laisse environ un dollar par boisson.
\end{culture}
\begin{culture}[Au Royaume-Uni]
L'addition s'appelle \en{the bill}, et le service est souvent déjà inclus (\en{service charge}) :
vérifiez avant d'ajouter 10 \%. Dans un pub, on commande et on paie au bar, sans pourboire.
\end{culture}
\end{paire}

\begin{culture}[À table, à l'américaine]
L'eau du robinet est servie gratuitement, avec des glaçons. Les portions sont grandes : on peut
demander à emporter ses restes (\en{Can I have a box, please?}). Et dans beaucoup de restaurants, le
café filtre est resservi sans supplément (\en{free refills}).
\end{culture}

%% ---------------------------------------------------------------------
\rubrique{pratique}{Exercices}

\exercice[lire,niveau=1]{Le menu}{Lisez le menu, puis répondez.}
\begin{center}\otMenu\end{center}
\begin{questions}
\item How much is the chicken salad? \hfill \trou[52mm]{It's fourteen dollars.}
\item How much is the burger and fries? \hfill \trou[52mm]{It's sixteen dollars.}
\item How much is the fish of the day? \hfill \trou[52mm]{It's nineteen dollars.}
\item How much are two coffees? \hfill \trou[52mm]{They're six dollars.}
\item Is the water free? \hfill \trou[52mm]{Yes, it is.}
\end{questions}

\exercice[ecrire,niveau=1]{\en{I like} ou \en{I'd like} ?}{Complétez.}
\begin{questions}
\item \trou[22mm]{I'd like} a glass of water, please.
\item \trou[22mm]{I like} jazz. It's my favorite music!
\item \trou[22mm]{I'd like} the chicken salad, please.
\item \trou[22mm]{I like} coffee, but I don't like tea.
\item \trou[22mm]{I'd like} a large latte to go, please.
\item \trou[22mm]{I like} French food very much.
\end{questions}

\exercice[lire,niveau=1]{Questions et réponses}{Associez chaque phrase (1 à 5) à la réponse qui
convient (a à e).}
\begin{anglais}
\noindent\begin{tabularx}{\linewidth}{@{}>{\otSemi\fontsize{8.6}{10}\selectfont\color{otNavy}}l@{\hspace{2.4mm}}X@{\hspace{8mm}}>{\otSemi\fontsize{8.6}{10}\selectfont\color{otNavy}}l@{\hspace{2.4mm}}X@{}}
1. & Would you like some tea?        & a. & To go, please.\\[1.4mm]
2. & For here or to go?              & b. & Yes, the burger, please.\\[1.4mm]
3. & Are you ready to order?         & c. & No, it isn't.\\[1.4mm]
4. & Is the tip included?            & d. & Seven fifty.\\[1.4mm]
5. & How much is that?               & e. & Yes, please.\\
\end{tabularx}
\end{anglais}\par\vspace{3mm}
\noindent{\fontsize{10}{12}\selectfont\otSemi Réponses :}\hspace{4mm}%
1~→~\trou[9mm]{e}\hspace{6mm}2~→~\trou[9mm]{a}\hspace{6mm}3~→~\trou[9mm]{b}\hspace{6mm}%
4~→~\trou[9mm]{c}\hspace{6mm}5~→~\trou[9mm]{d}

\exercice[ecrire,niveau=2]{Au café}{Complétez avec \en{can}, \en{I'd}, \en{would}, \en{have},
\en{here} ou \en{go}.}
\begin{anglais}
\noindent\begin{tabular}{@{}l@{\hspace{6mm}}l@{}}
\otNomDit{sam} & Hi! What \trou{can} I get you?\\[1.8mm]
\otNomDit{jack} & Can I \trou{have} a coffee, please?\\[1.8mm]
\otNomDit{sam} & Sure. For \trou{here} or to \trou{go}?\\[1.8mm]
\otNomDit{jack} & To go. And \trou{I'd} like a muffin, please.\\[1.8mm]
\otNomDit{sam} & \trou{Would} you like a bag?\\[1.8mm]
\otNomDit{jack} & No, thanks.\\
\end{tabular}
\end{anglais}

\exercice[ecrire,niveau=2]{Soyez poli !}{Ces demandes sont trop directes. Récrivez-les poliment, avec
\en{Can I have…?}, \en{Could I have…?} ou \en{I'd like…}}
\begin{questions}
\item Give me a coffee. \hfill \trou[76mm]{Can I have a coffee, please?}
\item I want the fish. \hfill \trou[76mm]{I'd like the fish, please.}
\item Give us the check. \hfill \trou[76mm]{Can we have the check, please?}
\item I want a glass of water. \hfill \trou[76mm]{Could I have a glass of water, please?}
\end{questions}

\exercice[lire,niveau=1]{Au restaurant}{Remettez le dialogue dans l'ordre : écrivez les lettres.}
\begin{anglais}
\begin{questions}[1][english]
\item[a.] I'd like the soup of the day, please.
\item[b.] Good evening! Are you ready to order?
\item[c.] Just a glass of water, please.
\item[d.] Very good. Would you like something to drink?
\item[e.] Yes, I am.
\end{questions}
\end{anglais}
\par\vspace{2mm}\noindent{\otSemi\fontsize{10}{12}\selectfont Ordre :}\hspace{4mm}%
\trou[9mm]{b}\hspace{3mm}\trou[9mm]{e}\hspace{3mm}\trou[9mm]{a}\hspace{3mm}\trou[9mm]{d}\hspace{3mm}\trou[9mm]{c}

\exercice[ecouter,ecrire,niveau=2,piste=11.10]{Les formes faibles}{Écoutez chaque demande, dite deux
fois à vitesse normale, et écrivez les mots qui manquent.}
\begin{questions}
\item \trou[30mm]{Can I have} a tea, please?
\item \trou[30mm]{Could I have} the check?
\item \trou[30mm]{Would you like} some water?
\item \trou[30mm]{I'd like} a salad, please.
\item \trou[30mm]{Can we have} two coffees?
\item \trou[30mm]{I'll have} the steak.
\end{questions}

\exercice[ecouter,ecrire,niveau=2,piste=11.11]{Les commandes}{Trois clients commandent au café de
Sam. Écoutez et complétez le tableau.}
\begin{center}\otTableau
\begin{tabular}{@{}l@{\hspace{10mm}}l@{\hspace{10mm}}l@{\hspace{10mm}}l@{}}
\otTH{Client} & \otTH{Boisson} & \otTH{Taille} & \otTH{Sur place ou à emporter}\\\midrule
{\otSemi 1 · Emma} & \enb{\trou[24mm]{tea}}    & \enb{\trou[18mm]{small}}  & \enb{\trou[22mm]{for here}}\\[1mm]
{\otSemi 2 · Jack} & \enb{\trou[24mm]{coffee}} & \enb{\trou[18mm]{large}}  & \enb{\trou[22mm]{to go}}\\[1mm]
{\otSemi 3 · Maya} & \enb{\trou[24mm]{latte}}  & \enb{\trou[18mm]{medium}} & \enb{\trou[22mm]{to go}}\\\bottomrule
\end{tabular}\end{center}

\exercice[ecouter,ecrire,niveau=2,piste=11.12]{Dictée à trous}{Au restaurant, le serveur prend la
commande d'Emma. Écoutez et complétez.}
\begin{anglais}
\noindent\begin{tabular}{@{}l@{\hspace{6mm}}l@{}}
\otNomDit{serveur} & Are you \trou{ready} to order?\\[1.8mm]
\otNomDit{emma}    & Yes. \trou{I'd} like the soup, please.\\[1.8mm]
\otNomDit{serveur} & \trou{Would} you like something to drink?\\[1.8mm]
\otNomDit{emma}    & Yes, \trou{please}. Can I have a \trou{glass} of orange juice?\\[1.8mm]
\otNomDit{serveur} & Of course. And for \trou{dessert}?\\[1.8mm]
\otNomDit{emma}    & No, \trou{thank} you.\\
\end{tabular}
\end{anglais}

\exercice[ecrire,niveau=2]{Traduisez}{Traduisez en anglais.}
\begin{questions}[1][french]
\item Je voudrais un thé, s'il vous plaît.\lignes{1}
\item Sur place ou à emporter ? — À emporter.\lignes{1}
\item Voulez-vous quelque chose à boire ? — Oui, volontiers.\lignes{1}
\item Pourrais-je avoir l'addition, s'il vous plaît ?\lignes{1}
\item Je prendrai le burger avec des frites.\lignes{1}
\item Le pourboire n'est pas compris.\lignes{1}
\item J'aime le café, mais je voudrais un thé aujourd'hui.\lignes{1}
\item C'est délicieux !\lignes{1}
\end{questions}

\exercice[ecrire,niveau=3]{Au \en{Brooklyn Bistro}}{Avec le menu de l'exercice 1, écrivez la scène
du restaurant (six à huit répliques) : le serveur, vous et un ami. Commandez, puis demandez
l'addition.}
\lignes{7}

\exercice[ecouter,parler,niveau=3,piste=11.13]{À vous de parler !}{Vous êtes au café de Sam. Répondez
\textbf{à voix haute} pendant chaque pause.}
\begin{center}\otTableau
\begin{tabular}{@{}>{\otSemi\fontsize{8.6}{10}\selectfont\color{otNavy}}l@{\hspace{3mm}}>{\enOn}p{70mm}<{\enOff}@{\hspace{6mm}}p{78mm}@{}}
 & \otTH{Sam dit\ldots} & \otTH{Vous\ldots}\\\midrule
1. & Hi there! What can I get you? & commandez une boisson, poliment.\\
2. & Small, medium or large? & choisissez une taille.\\
3. & For here or to go? & répondez.\\
4. & Would you like a muffin with that? & acceptez ou refusez poliment.\\
5. & That's six dollars, please. & donnez l'argent (« Voici. ») et remerciez.\\\bottomrule
\end{tabular}\end{center}

%% ---------------------------------------------------------------------
\rubrique{bilan}{Faites le point}

\begin{essentiel}
\begin{minipage}[t]{0.56\linewidth}
\otEssTitre{Les phrases clés}
\renewcommand{\arraystretch}{1.32}
\begin{tabular}{@{}>{\enOn\otSemi\fontsize{9.4}{11}\selectfont}l<{\enOff}@{\hspace{3mm}}>{\fontsize{9.2}{11}\selectfont\color{white!78!otNavy}}l@{}}
Can I have a latte, please?  & Je peux avoir un latte ?\\
I'd like the fish.           & Je voudrais le poisson.\\
For here or to go?           & Sur place ou à emporter ?\\
Would you like some tea?     & Voulez-vous du thé ?\\
Yes, please. / No, thanks.   & Oui, volontiers. / Non, merci.\\
Can we have the check?       & L'addition, s'il vous plaît.\\
\end{tabular}
\end{minipage}\hfill
\begin{minipage}[t]{0.39\linewidth}
\otEssTitre{Demander}
{\enb{\fontsize{9.2}{11.6}\selectfont Can I…? · I'd like… · I'll have… · Could I…?}\par}
\otEssTitre{Proposer}
{\enb{\fontsize{9.2}{11.6}\selectfont Would you like…?}\par
{\fontsize{9}{11.4}\selectfont\color{white!78!otNavy}\en{I like} (j'aime) ≠ \en{I'd like} (je voudrais)}\par}
\otEssTitre{Prononciation}
{\fontsize{9.2}{11.6}\selectfont \en{Can I…?} \phon{kən} · \en{Yes, you can.} \phon{kæn}\par}
\end{minipage}
\end{essentiel}

\begin{jesais}
\item commander une boisson au café, sur place ou à emporter ;
\item commander un repas au restaurant et demander l'addition ;
\item demander poliment avec \en{Can I…?}, \en{Could I…?}, \en{I'd like…} ;
\item proposer avec \en{Would you like…?} et répondre \en{Yes, please} / \en{No, thank you} ;
\item faire la différence entre \en{I like} et \en{I'd like} ;
\item calculer un pourboire à l'américaine.
\end{jesais}

\clearpage
\begin{testexpress}
\qcm[c]{« Je voudrais un café » se dit :}{\en{I like a coffee.}}{\en{I want a coffee.}}{\en{I'd like a coffee.}}
\qcm[a]{« Sur place ou à emporter ? » se dit :}{\en{For here or to go?}}{\en{Here or there?}}{\en{To eat or to take?}}
\qcm[b]{\en{Would you like some tea?} — Réponse polie :}{\en{Yes.}}{\en{Yes, please.}}{\en{Yes, I would like.}}
\qcm[a]{« L'addition, s'il vous plaît » (US) se dit :}{\en{Can we have the check, please?}}{\en{Give us the addition.}}{\en{The note, please.}}
\qcm[c]{La forme la plus polie est :}{\en{I want the fish.}}{\en{Give me the fish.}}{\en{Could I have the fish, please?}}
\qcm[b]{\en{a tip}, c'est :}{l'addition}{le pourboire}{le menu}
\qcm[a]{Aux États-Unis, le pourboire au restaurant est d'environ :}{15 à 20 \%}{2 à 5 \%}{rien, il est compris}
\qcm[c]{Dans \en{Can I have a coffee?}, \en{can} se prononce :}{\phon{kæn}, accentué}{\phon{kɑːn}}{\phon{kən}, très faible}
\qcm[b]{« Je prendrai le burger » (au restaurant) se dit :}{\en{I take the burger.}}{\en{I'll have the burger.}}{\en{I have the burger.}}
\qcm[a]{\en{fries}, ce sont :}{des frites}{des œufs frits}{des fruits}
\end{testexpress}

\noindent{\fontsize{9.6}{12}\selectfont\color{otMuted}Niveau visé (CECRL, niveau A1) : « peut
commander un repas ou une boisson et demander poliment ce dont il a besoin ». Les corrigés de
l'unité sont page~\pageref{cor-u11-1} ; les textes des pistes d'exercices,
page~\pageref{trans-11.10}.\par}
